\section{训练数据与经验内化}\label{app:train}
\subsection{训练语料的构建与整理}\label{app:corpus}
发布的语料由两部分组成。\emph{基础游戏}即生成的项目本身，以及产生每个项目的运行：1,998 个去重后的游戏（1,446 个 Godot、552 个 Phaser），其中 1,070 个附有生成轨迹，另有分布在三个生成批次中的 910 棵完整的逐次试验运行树。\emph{监督数据}是我们用于训练的内容，分三个阶段从这些运行中抽取，汇总于表~\ref{tab:corpus}：S1 游戏生成、S2 规划和 S3 改进。生成轨迹由 GPT-5.5 产生，改进轨迹由 GLM-5.3-Flash 产生。规划并非自由撰写：1,237 份从已完成游戏的文件树中蒸馏得到，870 份来自更早的一个规划集，仅当其提到的每个路径都存在于对应游戏中时才予以保留。

\begin{table}[htbp]
\caption{\textbf{发布的训练语料，按已发布的数据表计数。}一个\emph{决策行}是一次助手回合连同其之前的历史，因此决策行会重复上下文，其 token 计数并不相互独立；轨迹行是同样的录制内容未经切窗的形式，应以其统计 token。经验证的改进轮次是指经独立的改进后验证判定已实现既定目标的轮次；其余轮次作为负例发布，而不是丢弃。}
\label{tab:corpus}
\centering\small
\begin{tabular}{@{}llllrrr@{}}
\toprule
阶段 & 单位 & 教师模型 & & Godot & Phaser & 合计 \\
\midrule
S1 生成 & 轨迹 & GPT-5.5 & & 1,929 & 284 & 2,213 \\
S1 生成 & 决策行 & GPT-5.5 & & 66,312 & 28,097 & 94,409 \\
S2 规划 & 规划 & 蒸馏 & & 1,979 & 129 & 2,108 \\
S3 改进 & 轮次 & GLM-5.3-Flash & & 3,730 & 273 & 4,003 \\
\quad 其中经验证 & 轮次 & — & & 1,869 & 144 & 2,013 \\
S3 改进智能体 & 决策行 & GLM-5.3-Flash & & 111,302 & 9,930 & 121,232 \\
\quad 其中经验证 & 决策行 & — & & 48,341 & 5,098 & 53,439 \\
\midrule
被改进的游戏 & 游戏 & — & & 345 & 43 & 388 \\
不同任务 & 任务 & — & & 2,189 & 322 & 2,511 \\
训练行 & 行 & — & & 183,063 & 38,689 & 221,752 \\
\bottomrule
\end{tabular}
\end{table}

\paragraph{轨迹序列化。}每个生成会话被序列化为单条多轮智能体轨迹，包含任务规格说明、助手输出、工具调用和工具结果。训练损失只施加于助手生成的输出，避免对共享的交互上下文进行重复监督。规划样本取自智能体在编码之前的初始实现计划，而之后对同一计划的进度更新被丢弃。一个决策行保留一个固定前缀（系统提示、任务，以及存在时的规划），然后在 32k token 预算内放入尽可能多的最近完整决策，最后放入目标决策。

\paragraph{过滤掉什么。}仅当会话产出了可执行产物时才保留生成轨迹；判断依据是产物是否存在，而不是智能体 SDK 报告的成功状态，因为后者并不能可靠地表明是否产出了可运行的项目。空轨迹和重复任务被移除。仅当一份规划引用的每个文件路径都存在于完成的游戏中、且其引擎与简报相符时，才保留该规划。当改进轮次的会话失败或其差异为空时，该轮次被丢弃；且仅当经过独立的改进后验证时才标记为正例——Godot 的 3,730 个轮次中有 1,869 个，Phaser 的 273 次尝试中有 144 次。冗长的工具输出（例如构建日志）会被从中间截断，保留所执行的命令及其结论。整个项目的差异被压缩为改进实际修改的源代码片段，使改进样本的中位大小缩小约 48 倍。

\paragraph{评估污染。}任何基线、消融或评估运行的内容都不会进入语料，140 个不同的基准任务 id 在构建时即被排除。发布的数据表事后又针对这些 id 重新扫描：在全部 13 种配置、共 2,511 个任务 id 的并集上，重叠为零。

\paragraph{记录缺口被标记而非隐藏。}有一个智能体 SDK 保存的是工具结果的一行 UI 摘要，而不是模型实际收到的内容。受影响的行被保留，并带有显式标记（{\small\texttt{row\_has\_summary\_only}、\texttt{target\_follows\_summary\_only}、\linebreak\texttt{tool\_results\_source}、\texttt{trajectory\_incomplete}}），因此一个干净的训练子集只是一个过滤条件，而不是另一个发布版本。这影响全部 284 条 Phaser 生成轨迹（其中 20.2\% 的行的目标决策紧跟在摘要之后），以及两个 Godot 改进分支中的一个（31.7\% 的多轮改进行）；另一个 Godot 分支是完整的。

\paragraph{数据脱敏。}每一行、每个元数据列和每份文档都逐项重写。绝对机器路径变为 \texttt{/workspace/\allowbreak game}、\hspace{0pt}\texttt{/workspace/\allowbreak harness} 或 \texttt{<path>}（403,280 处重写）；用户、主机、账户和内部项目名称变为 \texttt{<redacted>}（13,610 处）；模型网关主机、内部镜像和集群名称被替换；电子邮件地址和非回环 IP 字面量被移除；捕获的目录列表中的属主列被重写；26,577 个配置字段被就地中和，因此数据模式保持不变。转储了构建主机进程表的工具输出被直接删除而不是清洗（26 份列表，379 行）。科学内容——游戏代码和素材、细则文本、逐需求的分数、模型标识符、推理强度、耗时——均保持不变。清洗器由另一个独立编写的扫描器进行验证，后者包含 47 种身份和凭据模式，以及一个要求移除所有预埋泄漏的预埋泄漏自测；二者报告的命中数均为零。

\subsection{微调配置}\label{app:ft}
我们使用参数高效的监督微调对 Qwen3.8-27B\citep{qwen2026}进行微调：在 4 比特量化的基础权重\citep{dettmers2023}上使用 LoRA\citep{hu2022lora}。每条智能体轨迹被视为一个长训练文档，使模型能够在原始开发上下文中学习规划、编码、工具使用和改进决策。除非另有说明，所有训练配置使用相同的超参数；只有训练语料的组成不同。

\paragraph{训练数据混合。}这三个变体并不是三种独立的混合，而是三个嵌套的语料：每个都在前一个的基础上加入一个阶段的监督，因此 $S1 \subset S2 \subset S3$，表~\ref{tab:ablation} 中任意两行之间的差异都可归因于新加入的那个阶段。两种引擎一起训练；表~\ref{tab:mixture} 给出了组成。

\begin{table}[htbp]
\caption{\textbf{表~\ref{tab:ablation} 中各变体的训练语料。}每个阶段都向前一个语料添加若干行；“累计”一列即该变体所用的训练数据。保留简报列表中的任务在每个阶段都被过滤掉。}
\label{tab:mixture}
\centering\small
\begin{tabular}{@{}llrrrr@{}}
\toprule
新增阶段 & 单位 & Godot & Phaser & 新增行数 & 累计 \\
\midrule
S1 生成 & 智能体轨迹 & 876 & 418 & 1,294 & 1,294 \\
S2 规划 & 简报 $\rightarrow$ 文件规划 & 216 & 92 & 308 & 1,602 \\
S3 改进 & 改进轮次 & 509 & 0 & 509 & 2,111 \\
\midrule
\textit{合计} & & 1,601 & 510 & 2,111 & \\
\bottomrule
\end{tabular}
\end{table}

\paragraph{长上下文训练。}智能体开发轨迹可能长达数万 token。因此，我们使用 57,344 token 的最大上下文长度，并使用 Ulysses 式序列并行\citep{jacobs2023}将每个序列分布到四块 GPU 上。我们使用梯度检查点和融合的交叉熵实现来降低激活值和输出头的显存消耗。

\paragraph{训练目标。}对每条轨迹，原始交互上下文被完整保留，而监督损失只施加于助手生成的输出。这使模型学习开发决策的序列，而不把工具观察或环境反馈当作预测目标。完整的 \rsig{} 模型在规划、生成和经独立验证的改进样本的并集上训练。

\subsection{训练数据消融}\label{app:ablation}
表~\ref{tab:ablation} 给出了第~5 节比较的每个变体（图~\ref{fig:global}c 绘制了 \textsc{+Gen} 和完整模型）。\textsc{Base} 和完整模型分别对应表~\ref{tab:main} 中的 Qwen3.8-27B 和 Qwen3.8-27B (SFT) 行。

\begin{table}[htbp]
\caption{\textbf{用于经验内化的微调配置。}这些是所发布适配器的设置；表~\ref{tab:ablation} 中的每个变体都以此训练，仅在表~\ref{tab:mixture} 的语料上有所不同。}
\label{tab:ftcfg}
\centering\small
\begin{tabular}{@{}ll@{}}
\toprule
配置项 & 取值 \\
\midrule
基础模型 & Qwen3.8-27B \\
适配方式 & LoRA\citep{hu2022lora} \\
LoRA 秩 $r$ & 16 \\
LoRA $\alpha$ & 32 \\
LoRA dropout & 0.05 \\
LoRA 目标模块 & 所有线性投影 \\
秩稳定 LoRA（rsLoRA） & 是 \\
基础权重精度 & 4 比特 NF4\citep{dettmers2023} \\
双重量化 & 是 \\
计算精度 & bfloat16 \\
最大上下文长度 & 57,344 token \\
训练轮数（epoch） & 1 \\
优化器 & AdamW\citep{loshchilov2019}，融合实现 \\
学习率 & $5\times10^{-5}$ \\
学习率调度 & 余弦 \\
预热比例 & 0.03 \\
权重衰减 & 0.01 \\
每设备批大小 & 1 \\
梯度累积 & 无 \\
序列并行度 & 4 \\
硬件 & 4$\times$ H100 80 GB \\
\bottomrule
\end{tabular}
\end{table}

\begin{table}[htbp]
\caption{\textbf{按训练变体划分的 Godot 单次生成结果。}每一行都基于 140 个任务。}
\label{tab:ablation}
\centering\small
\begin{tabular}{@{}lccccc@{}}
\toprule
变体 & 机制 & 深度 & 视觉 & 美术 & 总分 $\uparrow$ \\
\midrule
\textsc{Base} & 41.2 & 33.4 & 38.6 & 38.3 & 37.07 \\
\textsc{+Gen} & 49.0 & 41.4 & 42.0 & 38.0 & 41.45 \\
\textsc{+Gen +Plan} & 48.5 & 40.8 & 42.6 & 39.5 & 41.77 \\
\textsc{+Gen +Plan +Improve} & \textbf{56.1} & \textbf{47.5} & \textbf{49.8} & \textbf{44.8} & \textbf{48.22} \\
\bottomrule
\end{tabular}
\end{table}

\section{扩展评估与案例研究}\label{app:extended}
\subsection{与多智能体“生成–验证”系统的比较}\label{app:vibegame}
\paragraph{为何做这项比较。}本文并不声称智能体能够改进它自己生成的游戏——已经有若干系统做到了这一点。本文声称的是，要\emph{可靠地}做到这一点，需要一个控制器：一个能判断提升空间在哪里、保留最佳版本并适时停止的组件。我们能够运行的最接近的公开系统是 VibeGame\footnote{\url{https://github.com/tettethu/VibeGame}，Apache-2.0。}，这是一个由八个角色组成的智能体团队，其流程本身已包含验证和优化角色。如果这种结构的团队能够自行弥合差距，那么控制器就只是一个装饰。

\paragraph{设置。}使用 GameCraft-Bench 的四十个 Phaser 任务，VibeGame 代码不作修改，其全部八个角色都设为 GPT-5.5，每个任务一次试验，全自动、无人工干预。构建使用与本文所有其他 Phaser 分支相同的验证器、细则和评审打分。

\paragraph{结果说明了什么。}一个本身已经会验证和优化的智能体团队得分为 30.5，而在相同任务上单次生成为 46.0：它的循环并没有把轮次转化为质量。在我们共有的十二个任务上，它达到 23.3，而我们 Phaser 分支的冻结 $P_0$ 已经是 50.1，\rsig{} 将这十二个任务提升到 62.4。流程中存在验证和优化角色，并不是带来持续改进的原因。

\begin{table}[htbp]
\caption{\textbf{VibeGame 与单次生成及 \rsig{} 的对比。}左：VibeGame 的四十个任务与我们四十任务 Phaser 开发子集共有的十二个任务，因此每一列都是同样的十二个游戏。右：VibeGame 自己的四十个任务，与其发布版本在这些任务上引用的单次生成参考结果对比。}
\label{tab:vibegame}
\centering\small
\begin{tabular}{@{}lcc@{}}
\toprule
& 共有的 12 个任务 & VibeGame 自己的 40 个任务 \\
\midrule
VibeGame & 23.3 & 30.5 \\
OpenGame + GPT-5.5，单次生成 & 50.1 & 46.0 \\
\quad + Play2Code & 56.4 & — \\
\rowcolor{rowlight}\quad + \rsig & \textbf{62.4} & — \\
\bottomrule
\end{tabular}
\end{table}

\paragraph{公平地解读。}每个任务只有一次试验，而且我们对该系统的配置仅限于设置其角色模型；对于 VibeGame 在其他模型或更大预算下的表现，我们不作任何论断。

\subsection{定性演化案例研究}\label{app:cases}
我们跟踪三个游戏从代码模型生成的构建到我们最终交付的构建的全过程。每个案例以一条时间线开头：30 轮自主运行中每三轮的官方得分、全局质量监控器持有的版本、饱和停止结束运行的轮次，以及在坐标轴断开之后、循环本身没有运行的那些阶段——冠军不再被超越后写下的阶段简报，以及其后的开发阶段。时间线上的每个构建随后都有一张自己的卡片：该轮或该阶段发现了什么、改变了什么，以及据以评判它的画面；截图上的框是目测放置的，绿色表示某样东西已经到位，红色表示仍然有误。同一个游戏的每个构建都由同一脚本驱动，因此其面板可以相互比较。

\subsubsection{在退化中守住最佳版本}\label{app:case1}
\begin{figure}[h]
\centering
\fig{12}
\caption{\textit{Lawn Guardians}，从生成的游戏到我们交付的构建。得分带给出 30 轮自主运行中每三轮的官方得分、全局质量监控器持有的版本，以及饱和停止在第 21 轮（提前九轮）交付该版本的星标位置。轮次带为每一轮给出一个单元格——改进轮或美术轮——其上方是改进后的验证结论。坐标轴断开之后是循环没有运行的阶段：冠军经受住三个检查点后写下的阶段简报，以及其后的三次迭代；它们是开发阶段而不是循环轮次，因此没有轮次编号。下方的每个构建都在图~\ref{fig:lg1}–\ref{fig:lg2} 中展开为一张卡片。}
\label{fig:lg-timeline}
\end{figure}

开发并不在运行结束的地方结束。全局质量监控器的停止准则关注的是运行自身的得分，而不是游戏本身：一旦冠军经受住三个检查点，循环的证据就不再指向任何地方，而一份从外部写下的阶段简报提供了循环已经失去的方向。其后又有三个阶段——操控、动画，以及配有原创配乐的场景布置——每一个改变的都是游戏玩起来的感受，而不是它做了什么。六个构建在图~\ref{fig:lg1}–\ref{fig:lg2} 中各以一张卡片展示。

\begin{figure}[p]
\centering
\fig{13}
\caption{\textit{Lawn Guardians} 从生成的构建到阶段简报。每张卡片给出该轮或该阶段发现了什么、改变了什么，以及据以评判它的画面；截图上的框为目测放置，绿色表示某样东西已经到位，红色表示仍然有误，每个构建都由同一脚本驱动，因此场景可以相互比较。\textbf{基础游戏：}生成器把规则做对了、把呈现做错了——它附带的美术躺在磁盘上无人使用。\textbf{自主开发：}十二轮把这些美术接入并加上命中反馈，由此产生的冠军再未被超越。\textbf{阶段 1：}冠军被保持了三个检查点后，循环已经用尽了它自己的证据，而一份从外部写下的简报——由 Claude Opus 5 与人类各写一份——让阳光落在草坪上、让僵尸待在各自的行内。}
\label{fig:lg1}
\end{figure}

\begin{figure}[p]
\centering
\fig{14}
\caption{阶段简报之后的三个开发阶段。\textbf{迭代 1} 是一项任何棋盘静止画面都无法展示的修复，因此改为用同一脚本下的两组前后对比来记录：在拉伸的画布上，点击被变换了两次，使植物落在离光标一格远的地方，并且豌豆命中后仍留在草坪上。\textbf{迭代 2} 替换了标题画面和行走循环，是唯一一个静止画面确实能承载的步骤。\textbf{迭代 3} 是场景布置和重写的配乐；种子栏和边缘出现在画面中，而音频则是任何图都无法展示的那一半。}
\label{fig:lg2}
\end{figure}

在 \textit{Lawn Guardians}（图~\ref{fig:lg-timeline}）中，进展并非单调：一个美术轮（R6）和一个反馈轮（R18）各自降低了得分，而全局质量监控器两者都没有采纳。它一直持有 R12（86.1，起点为 61.7），直到饱和停止在第 21 轮交付它。这里也显示出两个局限：验证回放可能太短而看不到真正的修复（R26），而官方评审可能会奖励一个美术遮挡了游戏区域的构建（R27）。Play2Code 的峰值为 81.5，最终为 65.7。

\subsubsection{同一检查点上的两份简报}\label{app:case2}
\textit{Alley Brawlers} 是这样一个案例：循环依据自身证据换来了布景，却失去了演员。三十轮自主开发在打斗背后画出了一条灯火街道——这是整个运行中最大的单项变化——而在同样这些轮次中，格斗者不再能被正确渲染：在交付的构建中，运动中的身体是一团撕裂的涂抹。二者相互抵消，饱和停止交付的冠军得分与基础游戏完全相同。

循环无法得出的观察是：没有人能赢。两位阶段 1 导演都玩了这个构建并指出了这一点，而按模型简报进行的后续运行第一次让得分发生了变化。一个从同一检查点重启、完全不给简报的对照分支达到 60.0；两份简报分别达到 61.5 和 62.0。随后，在阶段 1 冠军上写下的第二份简报针对的是仍然让比赛无法进行的问题——两个身体站在同一空间里——换来的是一场可读的交锋，而不是更高的分数。

\begin{figure}[h]
\centering
\fig{15}
\caption{\textit{Alley Brawlers}，从生成的游戏到我们交付的构建。得分带、轮次带和画面同图~\ref{fig:lg-timeline}。饱和停止在第 18 轮结束自主运行并交付第 9 轮。坐标轴断开之后是循环没有运行的阶段：两份阶段简报（各由一位亲自玩过当时构建的导演写成）以及三个开发阶段。阶段 2 和各迭代都从阶段 1 冠军继续，因此它们是并列的备选，而不是一个序列。}
\label{fig:ab-timeline}
\end{figure}

\begin{figure}[p]
\centering
\fig{16}
\caption{\textit{Alley Brawlers} 从生成的构建到第一份简报。卡片、框和脚本同图~\ref{fig:lg1}。\textbf{基础游戏：}在下着方块雨的扁平条带上进行的一场能运行的比赛。\textbf{自主开发：}场地出现了，格斗者却坏了，交付的构建得分与基础游戏完全相同。\textbf{阶段 1：}一位打了两场比赛的导演要求公平的交锋，以及始终保持为一个完整人物的格斗者；得分在整个运行中第一次发生变化。}
\label{fig:ab1}
\end{figure}

\begin{figure}[p]
\centering
\fig{17}
\caption{第二份简报和各开发阶段，都从阶段 1 冠军继续。\textbf{阶段 2：}第二位导演指出间距是玩家无法行动的原因，交锋由此从双方都变得清晰可读。\textbf{迭代 1–2} 重建了前端——标题从一堵文字墙变成分步流程，HUD 为两名格斗者署名并展示其形象——并加入命中火花和伤害数字。\textbf{迭代 3} 关乎接触的手感：击中停顿、震屏、K.O. 慢动作、连击计数器和配乐。}
\label{fig:ab2}
\end{figure}

\subsubsection{一个循环无法改进的游戏}\label{app:case3}
\textit{Block Cascade} 是相反的起点。七袋随机器、后三块预览、幽灵方块、暂存交换、踢墙、消行、顶出判负和最高分在生成的构建中全都能正常工作，而三十轮由证据引导的改进找不到任何值得替换的版本。循环并不是卡在某个缺陷上；它已经找不到任何可以称为缺陷的东西了。

从该构建重新开始才带来了回报，而带来回报的并不是简报：两个由简报引导的分支和不给简报的对照分支都达到了 50.0，简报换来的是在第 6 轮而非第 12 轮到达那里。第二份简报随后要求增加紧张感——一个有名称的模式、一个实时计时器、一个奖励阶梯——并且得到了它们，同时也带来了新读数与屏幕上已有读数之间的冲突，这由其后的阶段重新布局解决。

\begin{figure}[h]
\centering
\fig{18}
\caption{\textit{Block Cascade}，从生成的游戏到我们交付的构建。生成的游戏已经实现了规格中的每个机制，因此改进轮几乎无事可改：在九个检查点中，全局质量监控器从未偏好过后来的构建，饱和停止在第 9 轮结束运行，而它交付的正是第 0 轮本身。坐标轴断开之后的一切都源自从该构建出发的重启。}
\label{fig:bc-timeline}
\end{figure}

\begin{figure}[p]
\centering
\fig{19}
\caption{\textit{Block Cascade} 从生成的构建经历两份简报。卡片、框和脚本同图~\ref{fig:lg1}。\textbf{基础游戏：}规格中的每个机制第一次就能工作，九轮之后循环交付的仍是这个构建。\textbf{阶段 1：}游戏井有了手绘背景、多面体方块和危险线——而一个不给简报的对照分支达到了同样的得分。\textbf{阶段 2：}有名称的模式和实时速度读数出现了，却被画在了已有读数之上。}
\label{fig:bc1}
\end{figure}

\begin{figure}[t]
\centering
\fig{20}
\caption{结束整个运行的两个开发阶段。\textbf{迭代 1} 为每个数字安排了独立的位置，并为落地加上反馈。\textbf{迭代 2} 增加了变暗棋盘上方的居中结算卡片、一个写明用哪个键填充的 HOLD 框，以及取代 1.8 秒片段的原创配乐。}
\label{fig:bc2}
\end{figure}

\subsection{按类别的结果}\label{app:family}
\newcommand{\famtab}[4]{%
\begin{table}[htbp]
\caption{#1}\label{#2}
\centering\small
\begin{tabular}{@{}lrccc@{}}
\toprule
类别 & $N$ & 基线 $\uparrow$ & #3 \\
\midrule
#4
\bottomrule
\end{tabular}
\end{table}}

\famtab{\textbf{Godot 上 Codex + GPT-5.5 (high) 生成器的按类别细分}（表~\ref{tab:main} 第一个行组），指标为平均总分（$\uparrow$），评审为 Qwen3.8-27B。$N$ 为该类别中的任务数（打分进行中时为“已打分/总数”）；各类别按基线得分排序；未生成项目的任务记 0 分。缩写：+P2C 即 Play2Code；–：尚未在全部任务上可得。}{tab:fam1}{+P2C & +\rsig}{
idle & 4 & 81.9 & 86.8 & 87.5 \\
sports & 4 & 72.8 & 70.5 & 81.2 \\
horror & 5 & 67.2 & 71.4 & 75.1 \\
openworld & 15 & 59.2 & 57.9 & 62.2 \\
tycoon & 16 & 53.3 & 54.4 & 73.2 \\
visualnovel & 11 & 52.0 & 55.8 & 71.5 \\
roguelike & 14 & 46.1 & 38.5 & 56.8 \\
shooter & 7 & 46.0 & 49.7 & 64.1 \\
strategy & 17 & 46.0 & 46.3 & 61.5 \\
rhythm & 5 & 45.4 & 44.7 & 67.9 \\
cardgame & 5 & 44.2 & 44.3 & 71.3 \\
puzzle & 8 & 43.4 & 45.1 & 53.9 \\
platformer & 19 & 43.3 & 44.4 & 58.4 \\
racing & 4 & 43.3 & 53.4 & 68.4 \\
simulation & 6 & 38.1 & 37.7 & 48.8 \\
\midrule
全部 & 140 & 50.26 & 50.74 & 64.53 \\}

\famtab{\textbf{Godot 上 Kimi-K2.6 生成器的按类别细分}（表~\ref{tab:main} 第二个行组）。}{tab:fam2}{+P2C & +\rsig}{
idle & 4 & 53.5 & 50.1 & 60.0 \\
sports & 4 & 45.9 & 57.3 & 64.4 \\
horror & 5 & 43.1 & 56.6 & 63.4 \\
tycoon & 16 & 36.1 & 40.7 & 46.4 \\
rhythm & 5 & 33.7 & 37.1 & 45.4 \\
shooter & 7 & 33.4 & 35.9 & 39.1 \\
simulation & 6 & 30.7 & 29.8 & 34.9 \\
visualnovel & 11 & 30.1 & 36.8 & 45.3 \\
roguelike & 14 & 29.4 & 35.1 & 44.5 \\
racing & 4 & 29.1 & 39.4 & 41.8 \\
openworld & 15 & 27.0 & 33.6 & 42.5 \\
cardgame & 5 & 26.2 & 23.6 & 48.8 \\
strategy & 17 & 24.1 & 33.1 & 47.7 \\
platformer & 19 & 23.3 & 29.2 & 41.4 \\
puzzle & 8 & 15.3 & 19.4 & 29.2 \\
\midrule
全部 & 140 & 29.63 & 35.19 & 44.77 \\}

\famtab{\textbf{Godot 上 GLM-5.3-Flash 生成器的按类别细分}（表~\ref{tab:main} 第三个行组）。}{tab:fam3}{+P2C & +\rsig}{
sports & 4 & 49.6 & 61.5 & 72.6 \\
idle & 4 & 39.0 & 59.2 & 68.6 \\
tycoon & 16 & 38.8 & 43.1 & 56.1 \\
racing & 4 & 34.9 & 36.7 & 53.0 \\
horror & 5 & 34.1 & 43.6 & 49.9 \\
visualnovel & 11 & 32.7 & 38.8 & 59.0 \\
shooter & 7 & 32.6 & 42.5 & 54.9 \\
strategy & 17 & 32.2 & 39.2 & 42.5 \\
simulation & 6 & 32.0 & 38.6 & 46.5 \\
rhythm & 5 & 26.3 & 37.2 & 52.9 \\
cardgame & 5 & 25.5 & 33.1 & 52.7 \\
platformer & 19 & 25.3 & 35.7 & 48.7 \\
roguelike & 14 & 24.9 & 37.5 & 40.9 \\
openworld & 15 & 19.2 & 28.1 & 43.2 \\
puzzle & 8 & 16.9 & 33.3 & 40.8 \\
\midrule
全部 & 140 & 29.46 & 38.59 & 49.72 \\}

\famtab{\textbf{Godot 上 Qwen3.8-27B 生成器的按类别细分}（表~\ref{tab:main} 第四个行组）。}{tab:fam4}{+P2C & +\rsig}{
horror & 5 & 58.7 & 56.6 & 57.0 \\
idle & 4 & 48.8 & 53.2 & 73.9 \\
cardgame & 5 & 48.3 & 53.1 & 59.3 \\
sports & 4 & 44.8 & 51.2 & 55.6 \\
shooter & 7 & 42.3 & 44.7 & 61.5 \\
visualnovel & 11 & 42.0 & 46.0 & 44.9 \\
tycoon & 16 & 41.0 & 52.0 & 62.6 \\
strategy & 17 & 39.5 & 35.2 & 44.9 \\
openworld & 15 & 34.5 & 39.7 & 36.4 \\
rhythm & 5 & 33.6 & 35.0 & 54.7 \\
simulation & 6 & 32.1 & 36.4 & 33.3 \\
platformer & 19 & 31.6 & 31.7 & 43.5 \\
roguelike & 14 & 31.0 & 36.7 & 42.4 \\
puzzle & 8 & 23.5 & 32.5 & 43.5 \\
racing & 4 & 23.3 & 22.7 & 27.9 \\
\midrule
全部 & 140 & 37.07 & 40.53 & 47.77 \\}

\famtab{\textbf{Godot 上 Qwen3.8-27B (SFT) 生成器的按类别细分}（表~\ref{tab:main} 第五个行组）。}{tab:fam5}{Play2Code $\uparrow$ & \rsig{} $\uparrow$}{
idle & 4 & 80.4 & 78.6 & 84.1 \\
sports & 4 & 73.3 & 68.7 & 76.1 \\
horror & 5 & 60.6 & 62.4 & 71.5 \\
tycoon & 16 & 54.7 & 56.4 & 66.3 \\
strategy & 17 & 52.1 & 51.0 & 62.0 \\
visualnovel & 11 & 50.2 & 48.9 & 58.4 \\
openworld & 15 & 48.1 & 46.7 & 58.2 \\
roguelike & 14 & 46.4 & 47.8 & 60.2 \\
platformer & 19 & 42.3 & 46.9 & 56.0 \\
racing & 4 & 41.3 & 38.5 & 52.1 \\
puzzle & 8 & 40.6 & 44.8 & 60.3 \\
simulation & 6 & 40.4 & 45.1 & 61.2 \\
rhythm & 5 & 37.4 & 43.9 & 63.8 \\
cardgame & 5 & 36.1 & 41.8 & 58.9 \\
shooter & 7 & 35.3 & 39.8 & 55.2 \\
\midrule
全部 & 140 & 48.22 & 49.71 & 61.38 \\}

\famtab{\textbf{Phaser 上 OpenGame + GPT-5.5 生成器的按类别细分}（表~\ref{tab:main} Phaser 部分第一个行组）。}{tab:fam6}{+P2C & +\rsig}{
idle & 4 & 72.7 & 73.8 & 78.2 \\
cardgame & 5 & 68.9 & 68.9 & 61.8 \\
visualnovel & 11 & 68.0 & 68.6 & 73.5 \\
tycoon & 16 & 59.9 & 62.9 & 74.3 \\
sports & 4 & 57.2 & 59.8 & 64.8 \\
roguelike & 14 & 56.5 & 58.2 & 66.7 \\
horror & 5 & 53.5 & 52.7 & 57.6 \\
simulation & 6 & 51.8 & 57.7 & 55.4 \\
rhythm & 5 & 50.8 & 52.5 & 58.2 \\
strategy & 17 & 45.7 & 56.8 & 54.8 \\
shooter & 7 & 45.2 & 58.6 & 60.4 \\
racing & 4 & 39.8 & 60.4 & 52.8 \\
platformer & 19 & 38.1 & 43.7 & 45.0 \\
openworld & 15 & 32.8 & 41.0 & 39.5 \\
puzzle & 8 & 32.5 & 40.6 & 51.6 \\
\midrule
全部 & 140 & 49.44 & 55.10 & 58.21 \\}

\famtab{\textbf{Phaser 上 Qwen3.8-27B 生成器的按类别细分}（表~\ref{tab:main} Phaser 部分第二个行组）。}{tab:fam7}{+P2C & +\rsig}{
sports & 4 & 56.1 & 68.8 & 65.4 \\
horror & 5 & 55.6 & 52.9 & 57.6 \\
simulation & 6 & 55.3 & 53.3 & 67.8 \\
visualnovel & 11 & 53.3 & 54.2 & 52.4 \\
rhythm & 5 & 50.9 & 58.4 & 61.4 \\
tycoon & 16 & 47.5 & 58.9 & 56.7 \\
cardgame & 5 & 47.0 & 72.3 & 61.4 \\
puzzle & 8 & 42.6 & 44.2 & 39.5 \\
idle & 4 & 41.3 & 61.5 & 65.2 \\
racing & 4 & 40.9 & 45.3 & 51.6 \\
strategy & 17 & 37.3 & 50.0 & 50.6 \\
platformer & 19 & 37.1 & 40.2 & 46.2 \\
roguelike & 14 & 29.6 & 40.7 & 48.8 \\
shooter & 7 & 27.7 & 46.0 & 48.1 \\
openworld & 15 & 21.1 & 32.0 & 29.3 \\
\midrule
全部 & 140 & 40.03 & 48.70 & 50.24 \\}

\famtab{\textbf{Phaser 上 Qwen3.8-27B (SFT) 生成器的按类别细分}（表~\ref{tab:main} Phaser 部分第三个行组）。}{tab:fam8}{Play2Code $\uparrow$ & \rsig{} $\uparrow$}{
idle & 4 & 73.0 & 75.2 & 78.5 \\
horror & 5 & 56.0 & 60.7 & 67.5 \\
tycoon & 16 & 53.3 & 62.0 & 66.5 \\
sports & 4 & 51.6 & 49.9 & 56.8 \\
shooter & 7 & 41.8 & 51.4 & 55.4 \\
rhythm & 5 & 46.8 & 61.3 & 66.2 \\
simulation & 6 & 46.0 & 50.7 & 60.1 \\
roguelike & 14 & 45.5 & 56.8 & 60.2 \\
openworld & 15 & 42.2 & 46.0 & 50.7 \\
puzzle & 8 & 43.8 & 58.7 & 64.8 \\
strategy & 17 & 42.9 & 51.3 & 54.9 \\
cardgame & 5 & 41.6 & 45.1 & 58.7 \\
visualnovel & 11 & 40.0 & 46.2 & 54.3 \\
platformer & 19 & 39.9 & 50.4 & 54.1 \\
racing & 4 & 32.8 & 39.1 & 49.6 \\
\midrule
全部 & 140 & 45.14 & 53.15 & 58.53 \\}
